\section{Introduction}
\label{sec:intro}

The bid-rigging detection literature has a blind spot. Its most
widely deployed tools---within-tender variance, kurtosis, and
coefficient of variation
\citep{imhof2019detecting, huber2019machine}---all rely on bid
dispersion. Yet the best-organized cartels defeat precisely these
screens: when cover bidders coordinate their bids tightly around
the designated winner, dispersion \emph{falls}, and the
statistical signature of collusion vanishes. The cartels that
matter most are the ones current tools cannot see.

We exploit a different footprint. Cover bidding requires firms to
show up and lose---repeatedly. A firm that bids on 50 tenders
over a decade without winning a single one is either
extraordinarily unlucky or playing a different game. We call these
firms \emph{frequent losers} (FL) and show that this participation
anomaly---observable from win/loss records alone, without any bid
microdata---marks procurement environments with systematically higher
prices and deep proximity to known cartels.

A cover-bidding model (Section~\ref{sec:structural}) formalizes the
intuition and distinguishes two regimes: under complementary cover
bidding (Regime~1), cover bids scatter above the winning bid; under
coordinated cover bidding (Regime~2), they cluster near a focal
point. The distinction has a sharp detection implication: Regime~2
defeats dispersion-based screens precisely when coordination is
strongest. A participation-based screen is immune to this problem
because it does not depend on how cover bids are distributed---only
on the fact that cover bidders must participate.

We test this logic on S\~ao Paulo's Bolsa Eletr\^onica de Compras (BEC),
covering 4.5 million tender-items and 40 million bids over 2009--2019.
Among 41,000 firms, 16,843 never win a single tender; a simple
IQR-based outlier rule identifies 2,735 whose participation frequency
cannot be squared with profit-maximizing entry---the frequent losers.
Against the same CADE ground truth, the rule achieves higher AUC
than a random forest and a structural classifier
(Section~\ref{sec:classification_robustness}), though the
comparison uses coarser proxies for the full Imhof implementation.
The single most discriminative feature of cover
bidding---participating heavily despite never winning---is hard for
more complex methods to improve upon.

Validated against CADE (Brazil's competition authority), the screen
achieves AUC~$= 0.94$, with the data-driven optimal threshold nearly
identical to the \emph{a priori} rule. FL firms co-participate with
convicted cartelists at 3.5$\times$ the baseline rate ($p < 0.001$).
Three CADE-convicted cover bidders are themselves classified as FL
(Section~\ref{sec:cade}).

FL-flagged tenders show 3.6--7.7\% higher prices in within-item,
within-year, within-purchasing-unit comparisons. The range spans
temporal cross-fitting (3.6\%), OLS with high-dimensional fixed
effects (6.4\%), and matching (CEM: 7.7\%; IPW: 5.5\%). After
absorbing item fixed effects the raw gap collapses from 2.43 to 0.050
log points (overlap 0.978; Table~\ref{tab:conditional_descstats},
Appendix~\ref{app:competition}), so the regressions compare
observationally similar tenders within tight product categories.
Sensitivity analysis (\citealt{cinelli2020making},
$RV = 17.5\%$) sets a high bar for omitted confounders.

The model generates five testable predictions; we take each to the
data. The price effect concentrates in competitive markets, where the
model says cover bidding pays the highest marginal return. BIC selects
Regime~2, revealing a \emph{dispersion paradox}: cover-bid dispersion
is 28\% \emph{below} genuine bids ($\sigma_c/\sigma_g = 0.72$),
the opposite of what variance-based screens assume.
Bajari--Ye tests reject exchangeability of FL and non-FL bid
residuals. FL--winner pairs recur far more often than chance
predicts. And the FL premium is larger in tenders where cover bidding
is voluntary (7.6\%) than where the minimum-bidder rule forces
participation. A horse-race regression shows FL captures price
information nearly orthogonal to \citet{imhof2019detecting}-style
bid-level features (correlation 0.06). No single test is dispositive,
but the joint pattern is difficult to generate under any non-collusive
story.

\subsection{Contributions}

\emph{First}, we develop a model of cover bidding that exposes a
blind spot in existing detection tools. The model distinguishes
complementary from coordinated regimes and shows that the coordinated
regime---BIC-selected in the data---defeats dispersion-based screens.
This makes participation-based detection not a convenience but a
necessity.

\emph{Second}, we build a screen that requires only win/loss
records---no bid microdata, no enforcement priors, no supervised
training labels. It achieves AUC~$= 0.94$ against
competition-authority convictions, achieving higher AUC than a
random forest and a structural classifier on the same ground
truth. The screen is
deployable wherever procurement systems record who bid and who won.

\emph{Third}, we document a 3.6--7.7\% price gap in FL-flagged
tenders and trace it through five independent diagnostics that jointly
match the coordinated cover-bidding model: regime selection,
Bajari--Ye bid-coordination tests, network-split heterogeneity,
dyadic linkage, and invitation variation. The challenge for
alternative explanations is to produce all five patterns
simultaneously.

We emphasize what the paper does \emph{not} claim. The price
association is a conditional correlation, not a causal estimate; the
FL classification identifies a participation anomaly, not true cartel
membership; and the results are specific to one procurement platform.
The paper's contribution is the screen and its diagnostic validation,
not causal identification of the FL--price effect
(Section~\ref{sec:limitations}).
